Une charge ponctuelle $Q = \qty{\v{Q}}{nC}$ est fixée à l'origine. Le point $A$ est à \qty{\v{rA}}{cm} de $Q$ et le point $B$ est à \qty{\v{rB}}{cm} de $Q$.

\begin{enumerate}[label=\alph*)]
    \item Déterminez le potentiel électrique aux points $A$ et $B$.
    \item On déplace une charge $q_0 = \qty{\v{q0}}{nC}$ de $A$ à $B$. Déterminez le travail effectué par la force électrique.
    \item Déterminez la variation d'énergie potentielle électrique de $q_0$. Le déplacement se fait-il spontanément ?
\end{enumerate}
